\section{Chilean school admissions}\label{sec:chile}
The Chilean school choice system records the courses a family ranks, the school's priority and lottery information, and the initial offer the family actually receives \citep{saedata,correaschool}. We can ask whether a hypothetical shortlist would have included that offer. To learn who would receive a place on the new lists, we must replay the allocation.

We fitted two ridge logistic offer predictors to the same 483,070 applicants in 2019, using eligible discretionary applicant--course edges as observations. Both use program demand, capacity, congestion, list length, grade, priority indicators and a smoothed program offer rate estimated from 2019 (leave-one-out for training). The second adds log rank, a first-choice indicator and a rank--congestion interaction. We call the first score rank-blind because it omits these features. Both pipelines use rank to truncate at the continuity entitlement and to order proposals. Both break equal fitted scores with the same seeded edge-specific random keys. The 2019 fit uses observed offers as training labels; no evaluation-year offer enters either policy. Evaluation-year demand is computed from the entire submitted list profile. A platform could use this score after reports are collected; advance advice would require a demand forecast. The lottery is used for allocation only. Each predictor retains up to five discretionary choices. A pupil's guaranteed current-school option remains available outside that quota. The eligible graph retains the report down to this entitlement, including the entitlement itself. In the sequel a full-list partner refers to DA on these eligible regular-round entries. The selected courses keep their submitted order when applicants propose. We use the same priority tiers, published within-course order and pooled capacities for both replays. For course $j$, the main seat count is $\max(v_j,h_j)$, where $v_j$ is the published vacancy count and $h_j$ the number of eligible continuing applicants. The latter adjustment makes every protected return feasible. Neither input uses evaluation-year offers. Full-list DA reproduces the recorded program or outside option for 95.38\% of applicants in 2022 and 95.78\% in 2023.

\paragraph{The first confirmation and its extension}
After examining the 2020 and 2021 application years, we recorded a joint prediction before obtaining the 2022 and 2023 releases. These application years correspond to school entry in 2023 and 2024. The cohorts include 520,420 and 492,888 initial offer recipients, respectively. All applicants enter the placement and joint-transition counts. Giving the 2019 predictor access to rank would cover more historical initial offers and reduce replayed placements in each subsequent cohort. Table~\ref{tab:chile-frozen} reports both outcomes, with full-list DA partners as a second, directly comparable coverage target. Among the 2,493 net additional historical offers covered in 2022--23, 2,244 were offers that families explicitly accepted. The historical record does not tell us whether they would have accepted under either new policy.

Before downloading either of the next two releases, we committed the five-slot comparison, the two directional criteria and the inclusion of both 2024 and 2025 application years. We applied the same 2019 coefficients, 2019 program frequencies, quota, capacity adjustments and priority ordering without refitting. The original models passed both criteria in each additional year. Across these two new cohorts of 937,405 application records the rank-informed predictor covered 3,059 more recorded offers and placed 802 fewer applicants. Among the 1,072 who lost a placement, 748 had identical personal menus under both policies. Full-list replay agreed with the recorded program or outside option for 96.59\% in 2024 and 96.51\% in 2025. We reproduced every pupil's assignment with two implementations of DA under both policies and full lists. The fixed-priority model uses the within-course order in each release. For school entry in 2026 (application year 2025), the ministry replaced a random tie-break with a deterministic applicant--school ordering; the published field still supplies the realized order \citep{newtie,saedata}.

\begin{table}[ht]\centering\small
\caption{Rank-informed minus rank-blind school menus. Placement and full-list DA coverage come from the replay. Historical offer coverage uses the recorded initial course. Each entry is a net count in the same application-year cohort.}
\label{tab:chile-frozen}
\begin{tabular}{@{}lrrrr@{}}\toprule
Year & Applicants & Historical offers shown & Full-list partners shown & Placements\\\midrule
2022 & 570,891 & $+1,143$ & $+1,261$ & $-468$\\
2023 & 536,353 & $+1,350$ & $+1,473$ & $-517$\\
2024 & 473,482 & $+1,735$ & $+1,841$ & $-481$\\
2025 & 463,923 & $+1,324$ & $+1,400$ & $-321$\\\bottomrule
\end{tabular}
\end{table}

The choice of comparison matters. In the 2022 and 2023 cohorts we also reran DA on 48 uniform five-choice displays in each cohort and on the family's first five choices. Relative to the uniform-draw mean, the rank-informed policy covered about 9,235 and 9,740 more full-list partners and placed about 789 and 759 more pupils in 2022 and 2023. Relative to first-five, it placed 2,123 and 2,013 more. The rank-blind predictor placed the most pupils of these four policies in both years. Prioritizing predicted offers improved on simple lists, and adding rank information then moved the two outcomes in opposite directions.


\paragraph{A market-wide predictor}
Across 2022--24, a rule using the completed applications to simulate unknown priority draws covered more historical offers and placed more pupils than the rank-informed predictor in every 64-draw batch. The placement advantage was small, from 2 to 56 pupils per batch, and each batch gave fewer first choices. When compared with the rank-blind rule it still placed fewer applicants. Appendix~\ref{sec:school-sensitivity} gives the information available to the simulated rule, its tie handling, and the results at smaller simulation budgets. It also reports seat, quota and priority-order sensitivities. Under a larger seat convention the replay fits the recorded assignment less well, although the original policy comparisons retain their signs.
